% Options for packages loaded elsewhere
\PassOptionsToPackage{unicode}{hyperref}
\PassOptionsToPackage{hyphens}{url}
%
\documentclass[
]{article}
\usepackage{amsmath,amssymb}
\usepackage{iftex}
\ifPDFTeX
  \usepackage[T1]{fontenc}
  \usepackage[utf8]{inputenc}
  \usepackage{textcomp} % provide euro and other symbols
\else % if luatex or xetex
  \usepackage{unicode-math} % this also loads fontspec
  \defaultfontfeatures{Scale=MatchLowercase}
  \defaultfontfeatures[\rmfamily]{Ligatures=TeX,Scale=1}
\fi
\usepackage{lmodern}
\ifPDFTeX\else
  % xetex/luatex font selection
\fi
% Use upquote if available, for straight quotes in verbatim environments
\IfFileExists{upquote.sty}{\usepackage{upquote}}{}
\IfFileExists{microtype.sty}{% use microtype if available
  \usepackage[]{microtype}
  \UseMicrotypeSet[protrusion]{basicmath} % disable protrusion for tt fonts
}{}
\makeatletter
\@ifundefined{KOMAClassName}{% if non-KOMA class
  \IfFileExists{parskip.sty}{%
    \usepackage{parskip}
  }{% else
    \setlength{\parindent}{0pt}
    \setlength{\parskip}{6pt plus 2pt minus 1pt}}
}{% if KOMA class
  \KOMAoptions{parskip=half}}
\makeatother
\usepackage{xcolor}
\usepackage{longtable,booktabs,array}
\usepackage{calc} % for calculating minipage widths
% Correct order of tables after \paragraph or \subparagraph
\usepackage{etoolbox}
\makeatletter
\patchcmd\longtable{\par}{\if@noskipsec\mbox{}\fi\par}{}{}
\makeatother
% Allow footnotes in longtable head/foot
\IfFileExists{footnotehyper.sty}{\usepackage{footnotehyper}}{\usepackage{footnote}}
\makesavenoteenv{longtable}
\setlength{\emergencystretch}{3em} % prevent overfull lines
\providecommand{\tightlist}{%
  \setlength{\itemsep}{0pt}\setlength{\parskip}{0pt}}
\setcounter{secnumdepth}{-\maxdimen} % remove section numbering
\ifLuaTeX
  \usepackage{selnolig}  % disable illegal ligatures
\fi
\IfFileExists{bookmark.sty}{\usepackage{bookmark}}{\usepackage{hyperref}}
\IfFileExists{xurl.sty}{\usepackage{xurl}}{} % add URL line breaks if available
\urlstyle{same}
\hypersetup{
  hidelinks,
  pdfcreator={LaTeX via pandoc}}

\author{}
\date{}

\begin{document}

\hypertarget{lezione-10-bitplanes}{%
\section{Lezione 10: BitPlanes}\label{lezione-10-bitplanes}}

Un'immagine con una profondità colore di N bit viene rappresentata da N piani di bit (bit-planes), ciascuno dei quali può essere vista come una singola immagine binaria. Si può indurre un ordine che varia dal Most Significant Bit (MSB) fino al Least Significant Bit (LSB).

💡

\textbf{Definizione:}

\textbf{Il bit plane di un'immagine digitale a N bit, è un'insieme di N immagini binarie (piani), in cui l'immagine i-esima contiene i valori dell'i-esimo bit della codifica scelta.}

Se tu \textbf{separi i bit}, ottieni \textbf{8 immagini diverse}.

Se si usa la codifica \textbf{binario puro} i piani di \textbf{bit più significativi} contengono informazioni sulla \textbf{struttura} dell'immagine, mentre quello via via \textbf{meno significativi} forniscono i \textbf{dettagli sempre più piccoli}.
Si nota che i piano dal 7 al 3 contengono dati significativi dal punto di vista visuale.
\textbf{Il rumore delle immagini e gli errori di acquisizione sono più evidenti nei piani bassi.}

Ogni bit-plane è:

\begin{itemize}
\tightlist
\item
  un'immagine \textbf{binaria}
\item
  pixel = 0 oppure 1
\item
  dice se \textbf{quel peso è presente o no}
\end{itemize}

Esempio:

Pixel = 127 → \texttt{01111111}

\begin{longtable}[]{@{}ll@{}}
\toprule\noalign{}
Piano & Valore \\
\midrule\noalign{}
\endhead
\bottomrule\noalign{}
\endlastfoot
piano 7 (128) & 0 \\
piano 6 (64) & 1 \\
piano 5 (32) & 1 \\
piano 4 (16) & 1 \\
piano 3 (8) & 1 \\
piano 2 (4) & 1 \\
piano 1 (2) & 1 \\
piano 0 (1) & 1 \\
\end{longtable}

\textbf{Ogni piano è un'immagine a sé}.

\begin{figure}
\centering
{image.png}
\caption{image.png}
\end{figure}

\begin{figure}
\centering
{image.png}
\caption{image.png}
\end{figure}

\begin{figure}
\centering
{Bit Plane di Lena}
\caption{Bit Plane di Lena}
\end{figure}

Bit Plane di Lena

\hypertarget{bit-planes-codifica-binario-puro}{%
\subsubsection{Bit planes: codifica binario puro}\label{bit-planes-codifica-binario-puro}}

Finora abbiamo detto che:
Un'immagine in scala di grigi \textbf{non è altro che una griglia di numeri}.
Ogni pixel ha un numero che dice \textbf{quanto è chiaro o scuro}.

\begin{itemize}
\tightlist
\item
  0 → nero
\item
  numero grande → più chiaro
\item
  massimo → bianco
\end{itemize}

Questi numeri devono essere \textbf{scritti da qualche parte nel computer}. Il computer però \textbf{non capisce i numeri decimali}, capisce solo \textbf{0 e 1}.

\begin{quote}
\textbf{Codifica in binario puro = modo standard per scrivere un numero usando solo 0 e 1.}
\end{quote}

``Binario puro'' significa:

\begin{itemize}
\tightlist
\item
  ogni numero viene scritto \textbf{in base 2}
\item
  ogni bit ha un \textbf{peso fisso}
\item
  quei pesi \textbf{non cambiano mai}
\end{itemize}

Poiché usiamo immagini ad 8 bit:

Questo significa:

\begin{itemize}
\item
  ogni pixel è un numero tra \textbf{\(0\) e \(255\)}
\item
  perché con 8 bit puoi fare:

  \(2^8=256\) valori → da \(0\) a \(255\)
\end{itemize}

I pesi dei bit sono:

\begin{longtable}[]{@{}ll@{}}
\toprule\noalign{}
Bit & Valore \\
\midrule\noalign{}
\endhead
\bottomrule\noalign{}
\endlastfoot
bit 7 & 128 \\
bit 6 & 64 \\
bit 5 & 32 \\
bit 4 & 16 \\
bit 3 & 8 \\
bit 2 & 4 \\
bit 1 & 2 \\
bit 0 & 1 \\
\end{longtable}

Si dice \textbf{binario puro} perché:

\begin{itemize}
\tightlist
\item
  ogni bit ha \textbf{sempre lo stesso significato}
\item
  il bit 7 vale \textbf{sempre 128}
\item
  il bit 0 vale \textbf{sempre 1}
\item
  non dipende dai pixel vicini
\end{itemize}

\textbf{Esempio chiave: pixel = 127}
Numero decimale: \textbf{127}
In binario puro: \texttt{01111111}

Traduzione:

\begin{itemize}
\tightlist
\item
  0×128 = 0
\item
  1×64 = 64
\item
  1×32 = 32
\item
  1×16 = 16
\item
  1×8 = 8
\item
  1×4 = 4
\item
  1×2 = 2
\item
  1×1 = 1
\end{itemize}

Somma: 64+32+16+8+4+2+1 = \textbf{127}
Ogni \textbf{1} significa ``quel contributo è acceso''.

\begin{figure}
\centering
{Codifica binario puro}
\caption{Codifica binario puro}
\end{figure}

Codifica binario puro

Questo genere di scomposizione è molto utile per eliminare tutti i valori compresi in un certo range.
Ad esempio se vogliamo eliminare tutti i grigi compresi tra 32 e 63, è necessari porre a 0 il quinto bit, e quindi tutto il piano 5.
Questa osservazione è valida solo se la codifica utilizzata è quella in binario puro.

\hypertarget{problema-del-bit-plane}{%
\subsubsection{Problema del Bit-Plane}\label{problema-del-bit-plane}}

Se la codifica usata è quella in \textbf{binario puro}, allora risulta evidente uno svantaggio: u\textbf{na piccola variazione può ripercuotersi su tutti i piani.}

Esempio: Se un pixel ha ad esempio intensità 127 (01111111) e il suo adiacente ha intensità 128 (10000000) allora la transizione tra 0 e 1 si ripercuote su tutti i piani di bit.

\textbf{Serve un codice in cui valori molti vicini abbiamo codifiche binarie molto simili, per evitare di incorrere in ``salti'' di intensità.}

\hypertarget{soluzione-codice-gray}{%
\subsubsection{Soluzione: Codice Gray}\label{soluzione-codice-gray}}

Il \textbf{codice Gray} è un modo alternativo di scrivere i numeri in binario tale che:

\begin{itemize}
\tightlist
\item
  ogni numero differisce dal precedente
\item
  \textbf{per un solo bit}
\item
  sempre
\end{itemize}

Non cambia il numero che rappresenti, cambia \textbf{come lo scrivi in binario}

\textbf{Attenzione: Gray NON è ``un altro numero''}

Questa è una cosa fondamentale.

\begin{itemize}
\tightlist
\item
  5 in decimale è \textbf{sempre 5}
\item
  ma:

  \begin{itemize}
  \tightlist
  \item
    in binario puro → \texttt{101}
  \item
    in Gray code → \texttt{111}
  \end{itemize}
\end{itemize}

Il valore è lo stesso.

La \textbf{codifica} è diversa.

💡

Il codice Gray a \(m\) bit \(g_{m-1} \dots g_1g_0\) che corrisponde al numero in binario puro \(a_{m-1} \dots a_1a_0\) può essere calcolato con la formula:

\[
g_i = a_i \oplus a_{i+1} \\ g_{m-1} = a_{m-1} \\
\text{con: } 0 \leq i \leq m - 2
\]

\begin{itemize}
\item
  \(a_i\) è il bit i-esimo del numero in binario puro, esempio:
  \(a_2 a_1 a_0 = 1 0 1\)
\item
  g\_i è il bit i-esimi del codice Gray
\item
  \(\oplus\) denota l'operatore XOR (OR esclusivo). Funziona così

  \begin{longtable}[]{@{}lll@{}}
  \toprule\noalign{}
  \(a\) & \(b\) & \(a \oplus b\) \\
  \midrule\noalign{}
  \endhead
  \bottomrule\noalign{}
  \endlastfoot
  0 & 0 & 0 \\
  0 & 1 & 1 \\
  1 & 0 & 1 \\
  1 & 1 & 0 \\
  \end{longtable}

  XOR vale \textbf{1 solo se i due bit sono diversi}
\end{itemize}

Ecco come calcolare il codice gray!

Ipotizziamo di voler trasformare un numero in 3 bit in codice gray.

\textbf{L'algoritmo è questo:}

\begin{enumerate}
\def\labelenumi{\arabic{enumi}.}
\tightlist
\item
  Converti in binario puro
\item
  Converti da binario puro in codice gray
\end{enumerate}

\textbf{Per convertire un numero DECIMALE in BINARIO con i pesi:}

\begin{enumerate}
\def\labelenumi{\arabic{enumi}.}
\item
  \textbf{Scrivi i pesi delle posizioni binarie}~da destra a sinistra:

  \begin{itemize}
  \tightlist
  \item
    In 3 bit: 4, 2, 1 (peso della posizione più a sinistra = 4, poi 2, poi 1)
  \end{itemize}
\item
  \textbf{Parti dal peso più grande}~(quello più a sinistra) e chiediti:

  ``Questo peso CI STA nel numero che sto convertendo?''
\item
  \textbf{Se CI STA}~(il peso è ≤ del numero rimanente):

  \begin{itemize}
  \tightlist
  \item
    Metti~\textbf{1}~in quella posizione
  \item
    \textbf{SOTTRAI}~il peso dal numero
  \end{itemize}
\item
  \textbf{Se NON CI STA}~(il peso è \textgreater{} del numero rimanente):

  \begin{itemize}
  \tightlist
  \item
    Metti~\textbf{0}~in quella posizione
  \item
    \textbf{NON sottrarre}~nulla
  \end{itemize}
\item
  \textbf{Passa al peso successivo}~(quello più piccolo) e ripeti
\end{enumerate}

\textbf{Per convertire un numero DECIMALE in BINARIO con la divisone per 2:}

\begin{enumerate}
\def\labelenumi{\arabic{enumi}.}
\tightlist
\item
  Dividi il numero per 2
\item
  Scrivi il~\textbf{RESTO}~(0 o 1)
\item
  Prendi il~\textbf{QUOZIENTE INTERO}~(senza decimali)
\item
  Ripeti finché il quoziente non diventa 0
\item
  \textbf{Leggi i resti DAL FONDO VERSO L'ALTO}~(dall'ultimo al primo)
\end{enumerate}

\textbf{Per convertire in codice gray:}

\begin{enumerate}
\def\labelenumi{\arabic{enumi}.}
\tightlist
\item
  \(g_2 = a_2\) \textbf{SEMPRE}
\item
  \(g_1 = a_2 \oplus a_1\)
\item
  \(g0 = a_1 \oplus a_0\)
\end{enumerate}

(Questo è nel caso dei 3 bit, in altre dimensioni i calcoli sono uguali, basta sostituire con i-esimo.)

\hypertarget{esempio-completo-da-decimale-a-binario-a-gray}{%
\subsubsection{Esempio Completo: Da Decimale a Binario a Gray}\label{esempio-completo-da-decimale-a-binario-a-gray}}

Facciamo un esempio \textbf{con il numero 5} e seguiamo TUTTI i passaggi che hai descritto.

\hypertarget{passo-1-conversione-decimale-binario-5-in-binario}{%
\subsubsection{\texorpdfstring{\textbf{PASSO 1: CONVERSIONE DECIMALE → BINARIO (5 in binario)}}{PASSO 1: CONVERSIONE DECIMALE → BINARIO (5 in binario)}}\label{passo-1-conversione-decimale-binario-5-in-binario}}

\hypertarget{metodo-a-divisioni-per-2}{%
\subsubsection{\texorpdfstring{\textbf{METODO A: DIVISIONI PER 2}}{METODO A: DIVISIONI PER 2}}\label{metodo-a-divisioni-per-2}}

\textbf{Numero da convertire: 5}

\begin{enumerate}
\def\labelenumi{\arabic{enumi}.}
\tightlist
\item
  \textbf{5 ÷ 2 = 2} con resto \textbf{1}

  \begin{itemize}
  \tightlist
  \item
    Quoziente: 2
  \item
    Resto: 1
  \end{itemize}
\item
  \textbf{2 ÷ 2 = 1} con resto \textbf{0}

  \begin{itemize}
  \tightlist
  \item
    Quoziente: 1
  \item
    Resto: 0
  \end{itemize}
\item
  \textbf{1 ÷ 2 = 0} con resto \textbf{1}

  \begin{itemize}
  \tightlist
  \item
    Quoziente: 0
  \item
    Resto: 1
  \end{itemize}
\end{enumerate}

\textbf{Ora abbiamo i resti in ordine:} 1, 0, 1

\textbf{Leggiamo DAL FONDO VERSO L'ALTO} (dall'ultimo resto al primo): 1, 0, 1

\textbf{Risultato:} \textbf{101₂} (binario a 3 bit)

\textbf{Verifica rapida:} 101₂ = (1×4) + (0×2) + (1×1) = 4 + 0 + 1 = 5 ✓

\hypertarget{metodo-b-pesi-4-2-1}{%
\subsubsection{\texorpdfstring{\textbf{METODO B: PESI (4, 2, 1)}}{METODO B: PESI (4, 2, 1)}}\label{metodo-b-pesi-4-2-1}}

\textbf{Numero da convertire: 5}

\textbf{Pesi per 3 bit:} 4, 2, 1

\begin{enumerate}
\def\labelenumi{\arabic{enumi}.}
\tightlist
\item
  \textbf{Peso 4} ci sta in 5? \textbf{SÌ} (4 ≤ 5)

  \begin{itemize}
  \tightlist
  \item
    Metti \textbf{1} nella posizione del peso 4
  \item
    5 - 4 = \textbf{1} (nuovo numero rimanente)
  \end{itemize}
\item
  \textbf{Peso 2} ci sta in 1? \textbf{NO} (2 \textgreater{} 1)

  \begin{itemize}
  \tightlist
  \item
    Metti \textbf{0} nella posizione del peso 2
  \item
    Numero rimanente resta \textbf{1}
  \end{itemize}
\item
  \textbf{Peso 1} ci sta in 1? \textbf{SÌ} (1 ≤ 1)

  \begin{itemize}
  \tightlist
  \item
    Metti \textbf{1} nella posizione del peso 1
  \item
    1 - 1 = \textbf{0}
  \end{itemize}
\end{enumerate}

\textbf{Risultato:} \textbf{101₂} (esattamente come con le divisioni!)

\textbf{Interpretazione:}

\begin{itemize}
\tightlist
\item
  Posizione del peso 4: \textbf{1}
\item
  Posizione del peso 2: \textbf{0}
\item
  Posizione del peso 1: \textbf{1}
\end{itemize}

\hypertarget{passo-2-da-binario-a-gray-101-binario-gray}{%
\subsubsection{\texorpdfstring{\textbf{PASSO 2: DA BINARIO A GRAY (101 binario → Gray)}}{PASSO 2: DA BINARIO A GRAY (101 binario → Gray)}}\label{passo-2-da-binario-a-gray-101-binario-gray}}

Ora applichiamo le tue formule per 3 bit:

\textbf{Binario:} \(( a_2 a_1 a_0 = 1 \ 0 \ 1 )\)
Dove:

\begin{itemize}
\tightlist
\item
  \(( a_2 = 1 )\) (bit più significativo, peso 4)
\item
  \(( a_1 = 0 )\) (bit centrale, peso 2)
\item
  \(( a_0 = 1 )\) (bit meno significativo, peso 1)
\end{itemize}

\textbf{Formule Gray:}

\begin{enumerate}
\def\labelenumi{\arabic{enumi}.}
\tightlist
\item
  \(( g_2 = a_2 = 1 )\) ← \textbf{SEMPRE si copia il primo bit}
\item
  \(( g_1 = a_2 \oplus a_1 = 1 \oplus 0 )\)
\item
  \((g_0 = a_1 \oplus a_0 = 0 \oplus 1\) )
\end{enumerate}

\hypertarget{calcoliamo}{%
\subsubsection{\texorpdfstring{\textbf{Calcoliamo:}}{Calcoliamo:}}\label{calcoliamo}}

\textbf{Per \(( g_1 ): ( 1 \oplus 0 )\)}

\begin{itemize}
\tightlist
\item
  XOR: 1 ⊕ 0 = \textbf{1} (perché 1 e 0 sono diversi)
\end{itemize}

\textbf{Per \(( g_0 ): ( 0 \oplus 1 )\)}

\begin{itemize}
\tightlist
\item
  XOR: 0 ⊕ 1 = \textbf{1} (perché 0 e 1 sono diversi)
\end{itemize}

\hypertarget{risultato-gray}{%
\subsubsection{\texorpdfstring{\textbf{Risultato Gray:}}{Risultato Gray:}}\label{risultato-gray}}

\begin{itemize}
\tightlist
\item
  \(( g_2 = 1 )\)
\item
  \(( g_1 = 1 )\)
\item
  \(( g_0 = 1)\)
\end{itemize}

\textbf{Codice Gray: 111}

\begin{center}\rule{0.5\linewidth}{0.5pt}\end{center}

\hypertarget{verifica-finale}{%
\subsection{\texorpdfstring{\textbf{VERIFICA FINALE}}{VERIFICA FINALE}}\label{verifica-finale}}

Abbiamo convertito: \textbf{5 (decimale) → 101 (binario) → 111 (Gray)}

\textbf{Verifichiamo con la tabella dei 3 bit:}

\begin{longtable}[]{@{}lll@{}}
\toprule\noalign{}
Decimale & Binario & Gray \\
\midrule\noalign{}
\endhead
\bottomrule\noalign{}
\endlastfoot
0 & 000 & 000 \\
1 & 001 & 001 \\
2 & 010 & 011 \\
3 & 011 & 010 \\
4 & 100 & 110 \\
\textbf{5} & \textbf{101} & \textbf{111} \\
6 & 110 & 101 \\
7 & 111 & 100 \\
\end{longtable}

\textbf{Perfetto!} 5 in binario è 101, e in Gray è 111.

\hypertarget{spiegazione-dettagliata-dellxor}{%
\subsection{\texorpdfstring{\textbf{SPIEGAZIONE DETTAGLIATA DELL'XOR (⊕)}}{SPIEGAZIONE DETTAGLIATA DELL'XOR (⊕)}}\label{spiegazione-dettagliata-dellxor}}

Forse ti stai chiedendo: ``Perché proprio queste formule per il Gray?''

L'operazione \textbf{XOR (⊕)} significa:

\begin{itemize}
\tightlist
\item
  \textbf{0 ⊕ 0 = 0} (stesso valore → 0)
\item
  \textbf{0 ⊕ 1 = 1} (valori diversi → 1)
\item
  \textbf{1 ⊕ 0 = 1} (valori diversi → 1)
\item
  \textbf{1 ⊕ 1 = 0} (stesso valore → 0)
\end{itemize}

In pratica: \textbf{Se i due bit sono UGUALI, risultato 0. Se sono DIVERSI, risultato 1.}

Nella conversione a Gray:

\begin{itemize}
\tightlist
\item
  \textbf{g₂ = a₂}: Il primo bit si copia sempre
\item
  \textbf{g₁ = a₂ ⊕ a₁}: Confronto tra il primo e il secondo bit del binario
\item
  \textbf{g₀ = a₁ ⊕ a₀}: Confronto tra il secondo e il terzo bit del binario
\end{itemize}

\hypertarget{secondo-esempio-completo-convertiamo-il-3}{%
\subsubsection{\texorpdfstring{\textbf{SECONDO ESEMPIO COMPLETO: CONVERTIAMO IL 3}}{SECONDO ESEMPIO COMPLETO: CONVERTIAMO IL 3}}\label{secondo-esempio-completo-convertiamo-il-3}}

\hypertarget{passo-1-3-in-binario}{%
\subsubsection{\texorpdfstring{\textbf{PASSO 1: 3 in binario}}{PASSO 1: 3 in binario}}\label{passo-1-3-in-binario}}

\textbf{Metodo divisioni:}

\begin{enumerate}
\def\labelenumi{\arabic{enumi}.}
\tightlist
\item
  3 ÷ 2 = 1 resto \textbf{1}
\item
  1 ÷ 2 = 0 resto \textbf{1}
  Resti: 1, 1 → Leggi al contrario: \textbf{11₂} → Ma servono 3 bit! Aggiungiamo uno 0 a sinistra: \textbf{011₂}
\end{enumerate}

\textbf{Metodo pesi (4,2,1):}

\begin{enumerate}
\def\labelenumi{\arabic{enumi}.}
\tightlist
\item
  Peso 4 in 3? NO → \textbf{0}
\item
  Peso 2 in 3? SÌ → \textbf{1}, 3-2=1
\item
  Peso 1 in 1? SÌ → \textbf{1}, 1-1=0
  Risultato: \textbf{011₂}
\end{enumerate}

\textbf{Binario: 011} (dove a₂=0, a₁=1, a₀=1)

\hypertarget{passo-2-da-binario-a-gray}{%
\subsubsection{\texorpdfstring{\textbf{PASSO 2: Da binario a Gray}}{PASSO 2: Da binario a Gray}}\label{passo-2-da-binario-a-gray}}

\begin{enumerate}
\def\labelenumi{\arabic{enumi}.}
\tightlist
\item
  \(( g_2 = a_2 = 0 )\)
\item
  \(( g_1 = a_2 ⊕ a_1 = 0 ⊕ 1 = 1 )\)
\item
  \(( g_0 = a_1 ⊕ a_0 = 1 ⊕ 1 = 0 )\)
\end{enumerate}

\textbf{Risultato Gray: 010}

\textbf{Verifica in tabella:} 3 in Gray è proprio 010 ✓

\begin{center}\rule{0.5\linewidth}{0.5pt}\end{center}

\hypertarget{terzo-esempio-convertiamo-il-6-per-esercizio}{%
\subsection{\texorpdfstring{\textbf{TERZO ESEMPIO: CONVERTIAMO IL 6 (per esercizio)}}{TERZO ESEMPIO: CONVERTIAMO IL 6 (per esercizio)}}\label{terzo-esempio-convertiamo-il-6-per-esercizio}}

Prova a farlo tu, poi controlla:

\textbf{Decimale 6 → Binario:}

\begin{itemize}
\tightlist
\item
  Metodo divisioni: 6÷2=3r0, 3÷2=1r1, 1÷2=0r1 → Resti: 0,1,1 → Al contrario: \textbf{110₂}
\item
  Metodo pesi: Peso 4 in 6? SÌ→1 (6-4=2), Peso 2 in 2? SÌ→1 (2-2=0), Peso 1 in 0? NO→0 → \textbf{110₂}
\end{itemize}

\textbf{Binario 110 → Gray:}

\begin{itemize}
\tightlist
\item
  a₂=1, a₁=1, a₀=0
\item
  g₂ = 1
\item
  g₁ = 1 ⊕ 1 = 0
\item
  g₀ = 1 ⊕ 0 = 1
\item
  \textbf{Risultato: 101}
\end{itemize}

Verifica in tabella: 6 in Gray è proprio 101 ✓

\hypertarget{riassunto-visivo-dellintero-processo-per-il-numero-5}{%
\subsubsection{\texorpdfstring{\textbf{RIASSUNTO VISIVO DELL'INTERO PROCESSO PER IL NUMERO 5}}{RIASSUNTO VISIVO DELL'INTERO PROCESSO PER IL NUMERO 5}}\label{riassunto-visivo-dellintero-processo-per-il-numero-5}}

\begin{verbatim}
DECIMALE: 5
        ↓
    CONVERSIONE A BINARIO
        ↓
METODO DIVISIONI:       METODO PESI:
5÷2=2r1                 Peso 4 in 5? SÌ → 1, 5-4=1
2÷2=1r0                 Peso 2 in 1? NO → 0
1÷2=0r1                 Peso 1 in 1? SÌ → 1
Resti: 1,0,1
Leggi al contrario: 101₂
        ↓
    BINARIO: 101 (a₂=1, a₁=0, a₀=1)
        ↓
    CONVERSIONE A GRAY
        ↓
Formule:
g₂ = a₂ = 1
g₁ = a₂ ⊕ a₁ = 1 ⊕ 0 = 1
g₀ = a₁ ⊕ a₀ = 0 ⊕ 1 = 1
        ↓
    GRAY: 111
\end{verbatim}

\hypertarget{lesempio-del-127-verso-128-il-caso-peggiore}{%
\subsubsection{L'esempio del ``127 verso 128'' (Il caso peggiore)}\label{lesempio-del-127-verso-128-il-caso-peggiore}}

Immagina un grigio medio (127) e un grigio appena appena più chiaro (128). L'occhio umano quasi non li distingue.

\begin{itemize}
\tightlist
\item
  \textbf{In Binario Puro:}
\end{itemize}

◦ 127 = \textbf{01111111}

◦ 128 = \textbf{10000000}

◦ \textbf{Disastro:} Sono cambiati \textbf{tutti e 8 i bit}! Se guardi i singoli piani bit-plane, vedresti un ``muro'' netto, come se l'immagine fosse cambiata drasticamente, mentre in realtà il colore è quasi uguale. Questo crea ``rumore'' visivo inutile nei piani bassi.

\begin{itemize}
\tightlist
\item
  \textbf{In Codice Gray:}
\end{itemize}

◦ 127 = \textbf{01000000}

◦ 128 = \textbf{11000000}

◦ \textbf{Soluzione:} Cambia \textbf{solo un bit}. La transizione è morbida anche a livello di dati.

Il Codice Gray ha una proprietà fondamentale: la \textbf{Distanza di Hamming (}distanza tra due stringhe, indica il numero di posizioni nelle quali i
simboli corrispondenti sono diversi
-numero di sostituzioni necessarie per convertire una stringa nell'altra\textbf{) è sempre 1}. Significa che tra un numero e il successivo cambia sempre e solo una cifra.

I vantaggi sono:

\begin{enumerate}
\def\labelenumi{\arabic{enumi}.}
\item
  \textbf{Coerenza tra i piani:} Se l'immagine ha colori sfumati (es. un cielo), i bit-plane in Codice Gray saranno più ``puliti'' e meno rumorosi rispetto al binario puro.
\item
  \textbf{Minore Entropia (più facile da comprimere):} Poiché ci sono meno cambiamenti frenetici di bit (meno transizioni bianco-nero inutili), il file contiene ``meno disordine''. Per un computer, un'immagine codificata in Gray è più facile da comprimere (occupa meno spazio) rispetto a una in binario puro, perché i dati sono più prevedibili.
\end{enumerate}

\textbf{In sintesi:} Se vuoi scomporre un'immagine in bit-plane per analizzarla o comprimerla, \textbf{non usare il binario classico}, perché crea ``falsi bordi'' (come tra 127 e 128). Usa il \textbf{Codice Gray}, che mantiene i dati stabili e coerenti con la percezione visiva.

\end{document}
